% Generisano: uređuj beleske/sNNN.tex i naslove u manifestu.
\documentclass[11pt,a4paper]{article}
% Zajednički preambulum za dokument za čitanje; nije Beamer tema.
\usepackage[a4paper,margin=20mm,headheight=15pt]{geometry}
\usepackage{lmodern,fontspec}
\usepackage[serbian]{babel}
\setmainfont{Latin Modern Sans}
\setsansfont{Latin Modern Sans}
\setmonofont{Latin Modern Mono}
\usepackage{amsmath,amssymb,mathtools,graphicx,xcolor}
\usepackage{array,booktabs,tabularx,multirow,makecell,listings,fancyhdr}
\pagestyle{fancy}\fancyhf{}
\newcommand{\BeleskeTekuciID}{NEOZNACEN}
\fancyhead[L]{\small Beleške uz slajd \texttt{\expandafter\detokenize\expandafter{\BeleskeTekuciID}}}
\fancyfoot[R]{\small\thepage}
\setlength{\parindent}{0pt}\setlength{\parskip}{6pt}
\newwrite\BeleskeMapa
\immediate\openout\BeleskeMapa=\jobname.notes-map.tsv
\AddToHook{shipout/before}{%
  \immediate\write\BeleskeMapa{\BeleskeTekuciID\space\thepage}}
\AddToHook{enddocument/afterlastpage}{\immediate\closeout\BeleskeMapa}
\newcommand{\BeleskeID}[1]{%
  \def\BeleskeZadatiID{#1}%
  \ifx\BeleskeTekuciID\BeleskeZadatiID\else
    \PackageError{beleske}{Pogresan ID beleske: #1}{Proveri mapiranje slajda.}%
  \fi}
\newcommand{\BeleskaSlajda}[4]{%
  \clearpage\gdef\BeleskeTekuciID{#1}%
  {\Large\bfseries Slajd \textnormal{\texttt{\detokenize{#1}}}: #3\par}\medskip
  \includegraphics[page=#2,width=\linewidth]{\BeleskePrezentacija}\par\medskip
  \input{#4}}

\newcommand{\BeleskePrezentacija}{build/07_uvod_u_hdl.pdf}
\begin{document}
\BeleskaSlajda{07-s001}{1}{Uvod u HDL}{beleske/s001.tex}
\BeleskaSlajda{07-s002}{2}{Opis digitalnog sistema}{beleske/s002.tex}
\BeleskaSlajda{07-s003}{3}{Fiksne i programabilne komponente}{beleske/s003.tex}
\BeleskaSlajda{07-s004}{4}{Ideja programabilne kombinacione mreže}{beleske/s004.tex}
\BeleskaSlajda{07-s005}{5}{Primer jedne programabilne komponente}{beleske/s005.tex}
\BeleskaSlajda{07-s006}{6}{Vrste programabilnih komponenti}{beleske/s006.tex}
\BeleskaSlajda{07-s007}{7}{Primer jedne PAL komponente}{beleske/s007.tex}
\BeleskaSlajda{07-s008}{8}{JEDEC}{beleske/s008.tex}
\BeleskaSlajda{07-s009}{9}{Razvoj HDL jezika}{beleske/s009.tex}
\BeleskaSlajda{07-s010}{10}{ABEL primeri}{beleske/s010.tex}
\BeleskaSlajda{07-s011}{11}{Primer}{beleske/s011.tex}
\BeleskaSlajda{07-s012}{12}{Tipičan CAD proces}{beleske/s012.tex}
\BeleskaSlajda{07-s013}{13}{Šta je VHDL?}{beleske/s013.tex}
\BeleskaSlajda{07-s014}{14}{VHDL: unos dizajna i testiranje}{beleske/s014.tex}
\BeleskaSlajda{07-s015}{15}{VHDL: netlist i standard}{beleske/s015.tex}
\BeleskaSlajda{07-s016}{16}{Kako se koristi VHDL?}{beleske/s016.tex}
\BeleskaSlajda{07-s017}{17}{VHDL: simulacija i dokumentacija}{beleske/s017.tex}
\BeleskaSlajda{07-s018}{18}{VHDL u procesu razvoja}{beleske/s018.tex}
\BeleskaSlajda{07-s019}{19}{VHDL i Verilog}{beleske/s019.tex}
\BeleskaSlajda{07-s020}{20}{Izbor HDL alata}{beleske/s020.tex}
\BeleskaSlajda{07-s021}{21}{Jednostavan VHDL kod}{beleske/s021.tex}
\BeleskaSlajda{07-s022}{22}{Opis kola – design entity - entity}{beleske/s022.tex}
\BeleskaSlajda{07-s023}{23}{Entity i architecture — obrazac}{beleske/s023.tex}
\BeleskaSlajda{07-s024}{24}{Jednobitni potpuni sabirač}{beleske/s024.tex}
\BeleskaSlajda{07-s025}{25}{„Dodeljivanje vrednosti“}{beleske/s025.tex}
\BeleskaSlajda{07-s026}{26}{Logičke i aritmetičke vektorske operacije}{beleske/s026.tex}
\BeleskaSlajda{07-s027}{27}{Prioritetni enkoder: ABEL i VHDL}{beleske/s027.tex}
\BeleskaSlajda{07-s028}{28}{Dataflow i redosled izraza}{beleske/s028.tex}
\BeleskaSlajda{07-s029}{29}{SR latch i delta ciklusi}{beleske/s029.tex}
\BeleskaSlajda{07-s030}{30}{Process — obrazac}{beleske/s030.tex}
\BeleskaSlajda{07-s031}{31}{MUX 2/1}{beleske/s031.tex}
\end{document}
